\label{chap:p3-83-23-algebra_clasificacion}


\label{int:chap:producto-nativo}
\index{native product}

The Euclidean residue--quotient extension of APP contains two local maps: the additive map normalizes collisions, and the multiplicative map produces partial products. In this section, the global multiplication of words is constructed using exclusively both maps. The correspondence with integer multiplication will be a consequence of the intertwiner, not the definition of the operation.

\section{Euclidean coordinates of the local cell}

\begin{definition}[Additive and multiplicative applications on the same support]
Let $\Dnine=\{1,\ldots,9\}$. For $a,b\in\Dnine$ they are defined by Euclidean division
with remainder in $\Dnine$:
\begin{align}
 a+b&=9q^+(a,b)+R^+(a,b),\label{int:eq:local-plus}\\
 ab&=9q^\times(a,b)+R^\times(a,b).\label{int:eq:local-times}
\end{align}
In particular, the zero remainder is represented by $9$ and the quotient is reduced by
one unit. The quadruple
\[
 \CellRQ(a,b)=
 (R^+(a,b),q^+(a,b);R^\times(a,b),q^\times(a,b))
\]
is the local arithmetic extension by residue and quotient.
\end{definition}

Aggregation by the three classes modulo three of the visible table
\(R^\times\) defines a matrix distinct from the local maps:
\[
(M_\Pi)_{rs}
=\frac19
\sum_{\substack{a\equiv r\;(3)\\b\equiv s\;(3)}}R^\times(a,b),
\qquad r,s\in\{1,2,0\},
\]
and, in that order of classes,
\[
M_\Pi=
\begin{pmatrix}4&5&6\\5&4&6\\6&6&9\end{pmatrix}.
\]
Thus, \(q^+\) and \(q^\times\) are quotient coordinates of a cell,
whereas \(M_\Pi\) is an aggregate operator on residual classes.

The $81$ cells satisfy exactly
\[
\begin{array}{c|rrrr}
 &\sum R&\sum q&\sum(9q+R)\\\hline
 +&405&45&810\\
 \times&459&174&2025
\end{array}
\]
and therefore
\[
 2025-810=(459-405)+9(174-45)=54+9\cdot129=1215.
\]
The quotient $q^\times$ is necessary for reconstruction. For example,
\[
 2\cdot5=9\cdot1+1.
\]
If only $R^\times$ is preserved, the product $10$ is conflated with the
visible value $1$.

\subsection{Cocycles of associativity and distributivity}

The coherence of the transducer can already be decided in the cell. For every
$a,b,c\in\Dnine$ the three identities are verified
\begin{align}
q^+(a,b)+q^+(R^+(a,b),c)
&=q^+(b,c)+q^+(a,R^+(b,c)),\label{int:eq:cocycle-plus}\\
q^\times(R^\times(a,b),c)+cq^\times(a,b)
&=q^\times(a,R^\times(b,c))+aq^\times(b,c),
\label{int:eq:cocycle-times}\\
q^\times(a,R^+(b,c))+aq^+(b,c)
&=q^+(R^\times(a,b),R^\times(a,c))\notag\\
&\quad+q^\times(a,b)+q^\times(a,c).
\label{int:eq:cocycle-distributive}
\end{align}

\begin{theorem}[Local coherence of the Euclidean extension]
The identities \eqref{int:eq:cocycle-plus}--\eqref{int:eq:cocycle-distributive},
together with the corresponding equalities of residues, express
respectively additive associativity, multiplicative associativity, and
distributivity in the residue--quotient extension. The three families have
zero discrepancies in the $9^3=729$ triples.
\end{theorem}

\begin{proof}
Expanding $(a+b)+c$ through \eqref{int:eq:local-plus}, the coefficient of
nine is the left-hand side of \eqref{int:eq:cocycle-plus}; expanding
$a+(b+c)$ yields the right-hand side. Uniqueness of division with remainder in
$\Dnine$ equates both residues and quotients. The same argument applied to
$(ab)c=a(bc)$ gives \eqref{int:eq:cocycle-times}. Finally, the two expansions of
$a(b+c)=ab+ac$ produce \eqref{int:eq:cocycle-distributive}. Direct enumeration
of the \(729\) cases of each identity confirms the algebraic
calculation.
\end{proof}

\section{Bijectivity of numbering}

The symbol $9$ is not identified with zero. In $\CellRQ$, $9$ represents a
complete crossing of the nonadic boundary. Replacing it with the digits
$0,\ldots,8$ would eliminate the Euclidean quotient that intervenes in the
intertwining.

\begin{definition}[Canonical words]
Let
\[
 \Wnine^+=\bigsqcup_{\ell\geq1}\Dnine^\ell,
 \qquad \Wnine=\{\varnothing\}\sqcup\Wnine^+.
\]
Words are written from the least significant digit. We define
\[
 \operatorname{ev}_9(u_0,\ldots,u_{\ell-1})
 =\sum_{j=0}^{\ell-1}u_j9^j,
 \qquad \operatorname{ev}_9(\varnothing)=0.
\]
\end{definition}

\begin{theorem}[Uniqueness of the normal form]
The evaluation $\operatorname{ev}_9$ is a bijection from $\Wnine^+$ onto
$\N_{\geq1}$, and from $\Wnine$ onto $\N_0$.
\end{theorem}

\begin{proof}
The words of length $\ell$ cubren the intervalo
\[
 \frac{9^\ell-1}{8}
 \leq \operatorname{ev}_9(u)
 \leq 9\frac{9^\ell-1}{8}.
\]
The minimum of the next interval is one unit greater than the maximum of the
preceding one. Within each interval, division of $n-1$ by nine uniquely determines
the digit
\[
 d=(n-1)\bmod9+1
\]
and the quotient that is encoded recursively. The process terminates because the
quotient strictly decreases.
\end{proof}

Examples:
\[
 9\longleftrightarrow(9),\qquad
 10\longleftrightarrow(1,1),\qquad
 81\longleftrightarrow(9,8),\qquad
 1000\longleftrightarrow(1,3,3,1).
\]

\section{Fibers of the residual projections on the integral liftings}

For this block, we distinguish two alphabets. The normal form
\(\Wnine\) uses the digits \(\{1,\ldots,9\}\), written from the least
significant. By contrast, \(\mathscr B_9\) denotes the conventional positional
representations in base nine, with digits \(\{0,\ldots,8\}\) and the most
significant on the left. The two types are not identified.

Let
\[
 \mathscr W_{10}^{(2)}=\{0,\ldots,9\}^2,\qquad
 \rho_{10}(a,b)=(b,a),\qquad
 \operatorname{ev}_{10}(a,b)=10a+b,
\]
and let \(q_9:\mathbb N\to\mathbb Z/9\mathbb Z\) be the residual reduction.

\begin{lemma}[Residual descent of decimal reversal]
\label{int:pf84:SYN30-RHO-DESC}
Reversion on decimal words of length two is involutive and its
shadow modulo nine is the identity:
\[
 \rho_{10}^2=I,\qquad
 q_9\!\left(\operatorname{ev}_{10}(\rho_{10}(a,b))\right)
 =
 q_9\!\left(\operatorname{ev}_{10}(a,b)\right).
\]
Fixed length preserves the initial zeros.
\end{lemma}

\begin{proof}
The doble permutation recovers \((a,b)\). Ifnce \(10\equiv1\pmod9\),
\[
 10b+a\equiv a+b\equiv10a+b\pmod9.
\]
A word whose double reversal were not itself, or whose digit sum
changed modulo nine, would refute the lemma.
\end{proof}

\begin{lemma}[Residual Square Descent]
\label{int:pf84:SYN30-SQ-DESC}
The map \(s(n)=n^2\) induces the residual endomap
\[
 [n]_9\longmapsto[n]_9^2,
 \qquad q_9(s(n))=q_9(n)^2.
\]
\end{lemma}

\begin{proof}
The compatibility of a congruence with the product gives
\(n\equiv m\pmod9\Rightarrow n^2\equiv m^2\pmod9\). An integer for which
\(q_9(n^2)\ne q_9(n)^2\) would refute the statement.
\end{proof}

\begin{proposition}[Residual fibers of the full reversal]
\label{int:pf84:SYN30-RHO-NONFAITH}
The residual shadow does not reconstruct the integral lift
\(\widehat\rho_{10}=\operatorname{ev}_{10}\circ\rho_{10}\). Indeed,
\[
 18\equiv27\pmod9,\qquad
 \widehat\rho_{10}(1,8)=81\ne72=\widehat\rho_{10}(2,7).
\]
\end{proposition}

\begin{proof}
Both integers belong to residue class zero, but their reversals are
distinct. Therefore, there is no
\(\overline\rho:\mathbb Z/9\mathbb Z\to\mathbb N\) that simultaneously factors
those two values.
\end{proof}

\begin{proposition}[Residual fibres of the entire square]
\label{int:pf84:SYN30-SQ-NONFAITH}
The same shadow likewise does not reconstruct the integer square:
\[
 18\equiv27\pmod9,\qquad
 18^2=324\ne729=27^2.
\]
\end{proposition}

\begin{proof}
A function \(\overline s:\mathbb Z/9\mathbb Z\to\mathbb N\) that
factorized \(s\) would have to assign the two values
\(324\) and \(729\) to the same zero class, which is impossible.
\end{proof}

\begin{corollary}[Cross-witness \(27\)--\(72\)--\(729\)]
\label{int:pf84:SYN30-CROSS-WITNESS}
\[
 \widehat\rho_{10}(2,7)=72,\qquad
 27^2=729=9^3=3^6,
\]
and, in the alfabeto convencional \(\mathscr B_9\),
\[
 27=30_9,\qquad72=80_9,\qquad729=1000_9.
\]
The three integers have residual shadow zero and positive digital root nine.
\end{corollary}

\begin{proof}
The two operations are computed directly; successive divisions by
nine produce the three expressions. The result would fail if
\(\mathscr B_9\) were confused with \(\Wnine\), if any of the expressions were incorrect,
or if it were claimed that faithful residual descent reconstructs the operation and its
complete genealogy.
\end{proof}

\section{Native adder and quotient propagation}

\begin{algorithm}[Word addition]
Two words are traversed by positions, with initial carry \(c_0=0\). If there are
two digits at a position, \eqref{int:eq:local-plus} is applied; if there is only one,
that digit is taken as the provisional residue and the provisional quotient is zero.
When \(c_j\ne0\), a second additive cell is applied to the provisional residue
and to \(c_j\). The new residue is published, and the sum of the
quotients is transported to the next position. If a nonzero carry remains upon
completion, it is published as the final digit. In this manner, the local cells never receive the
absent symbol \(0\), which does not belong to \(\Dnine\).
\end{algorithm}

\begin{lemma}[Adder Invariant]
If $u\boxplus_\#v$ is the word produced by the preceding algorithm,
\[
 \operatorname{ev}_9(u\boxplus_\#v)
 =\operatorname{ev}_9(u)+\operatorname{ev}_9(v).
\]
The intermediate carry lies in $\{0,1,2\}$.
\end{lemma}

\begin{proof}
After processing the positions $0,\ldots,k-1$, let $c_k$ be the pending quotient.
The local identities imply
\[
 \sum_{j<k}(u_j+v_j)9^j
 =\sum_{j<k}w_j9^j+c_k9^k,
\]
including the absent digits as zero. Each new cell preserves this
equality. The local maximum is $9+9+2=20$, so that $c_{k+1}\leq2$. At the
end the last quotient is published, and the global identity remains.
\end{proof}

\section{Multiplication by a digit}

\begin{algorithm}[Scalar transducer]
With $d\in\Dnine$ fixed, each digit $u_j$ produces
\[
 u_jd=9q_j+r_j,
 \qquad
 (r_j,q_j)=\bigl(R^\times(u_j,d),q^\times(u_j,d)\bigr).
\]
Let $c_j\in\{0,\ldots,9\}$ be the incoming quotient. Since $0\notin\Dnine$,
the transition is defined by branches:
\[
 (w_j,c_{j+1})=
 \begin{cases}
 (r_j,q_j),&c_j=0,\\[1mm]
 \bigl(R^+(r_j,c_j),q_j+q^+(r_j,c_j)\bigr),&c_j\in\Dnine.
 \end{cases}
\]
Thus, no local map is evaluated outside its domain. The first
branch expresses the absence of a pending quotient; the second normalizes the
collision by means of the additive map.
\end{algorithm}

\begin{lemma}[Scalar invariant]
The word $d\odot_\#u$ produced by the transducer satisfies
\[
 \operatorname{ev}_9(d\odot_\#u)
 =d\operatorname{ev}_9(u).
\]
The multiplicative quotient belongs to $\{0,\ldots,9\}$.
\end{lemma}

\begin{proof}
The argument is that of the adder, replacing each local summand by
$u_jd$. If $c_j=0$, the identity
$u_jd=9q_j+r_j$ directly preserves the partial equality. If
$c_j\in\Dnine$, the second Euclidean division gives
\[
 r_j+c_j=9q^+(r_j,c_j)+R^+(r_j,c_j),
\]
and preserves the same equality after transporting
$q_j+q^+(r_j,c_j)$. The maximum $9\cdot9+9=90$ and the remainder form in
$\Dnine$ keep the quotient within the declared bound. The last quotient is
published as a digit.
\end{proof}

\section{Global product by Horner}

\begin{definition}[Native product]
Let $v=(v_0,\ldots,v_{m-1})$. Partiendo of $a_m=\varnothing$, is defined for
$j=m-1,\ldots,0$
\[
 a_j=(9\odot_\#a_{j+1})\boxplus_\#(v_j\odot_\#u).
\]
The result is
\[
 u\boxtimes_\#v:=a_0.
\]
\end{definition}

This definition does not contain
\[
 \operatorname{dec}_9(\operatorname{ev}_9(u)\operatorname{ev}_9(v)).
\]
The evaluation will appear only in the proof of entanglement.

\begin{theorem}[Native multiplicativity]
For all words $u,v\in\Wnine$,
\begin{equation}
 \label{int:eq:ev-multiplicative}
 \boxed{
 \operatorname{ev}_9(u\boxtimes_\#v)
 =\operatorname{ev}_9(u)\operatorname{ev}_9(v).}
\end{equation}
The empty word is absorbing, and $(1)$ is the unit.
\end{theorem}

\begin{proof}
By the two previous lemmas,
\[
 \operatorname{ev}_9(a_j)
 =9\operatorname{ev}_9(a_{j+1})
  +v_j\operatorname{ev}_9(u).
\]
Horner's iteration gives
\[
 \operatorname{ev}_9(a_0)
 =\operatorname{ev}_9(u)
  \sum_{j=0}^{m-1}v_j9^j.
\]
The last factor is $\operatorname{ev}_9(v)$. The uniqueness of the normal form
identifies the resulting word.
\end{proof}

\begin{corollary}[Multiplicative monoid of words]
The structure
\[
\WordMonoid=(\Wnine,\boxtimes_\#,(1))
\]
is a commutative monoid; the empty word is its absorbing element. Its
construction does not use global multiplication of integers.
\end{corollary}

The shortest boundary example is
\[
 (9)\boxtimes_\#(9)=(9,8),
 \qquad 9+8\cdot9=81.
\]
The \(17\) local cells of multiplication produce
\[
 (9,8)\boxtimes_\#(9,8)=(9,8,8,8),
 \qquad \operatorname{ev}_9(9,8,8,8)=6561.
\]

\section{Matrix linearization of the genealogical state and construction of the completely positive entangler}

\begin{definition}[Hilbert Sectors]
Let
\[
 \Krad=\ell^2(\Wnine^+),
 \qquad
 \Hplus=\ell^2(\N_{\geq1}).
\]
We define on the canonical bases
\[
 W_\times e_u=e_{\operatorname{ev}_9(u)}.
\]
In \(\Hplus\), the number operator is denoted by
\[
 \mathsf N e_n=n e_n,
 \qquad \Spec(\mathsf N)=\N_{\geq1}.
\]
The vacuum may be added separately as a vacuum vector $e_0$; it is not
confused with $\Hplus$, whose spectrum begins at one.
\end{definition}

\begin{theorem}[Unitarity and multiplicative square]
$W_\times$ is unitary. In the algebraic kernel
$\mathbb C[\Wnine^+\times\Wnine^+]$, let
\[
 \mu_\#(e_u\otimes e_v)=e_{u\boxtimes_\#v},
 \qquad
 \mu_{\N}(e_a\otimes e_b)=e_{ab}.
\]
Then the square holds.
\[
\begin{CD}
\Krad\otimes_{\mathrm{alg}}\Krad @>{W_\times\otimes W_\times}>>
\Hplus\otimes_{\mathrm{alg}}\Hplus\\
@V{\mu_\#}VV @VV{\mu_{\N}}V\\
\Krad @>>{W_\times}> \Hplus
\end{CD}
\]
and therefore, the square commutes.
\end{theorem}

\begin{proof}
The bijection of normal forms sends one orthonormal basis to another and proves
unitarity. On a basic tensor, both paths produce
$e_{\operatorname{ev}_9(u)\operatorname{ev}_9(v)}$ by
\eqref{int:eq:ev-multiplicative}. Linearity gives the equality on the kernel.
\end{proof}

The maps $\mu_\#$ and $\mu_{\N}$, initially defined on the respective
algebraic kernels, are densely defined and unbounded. For
\(c=(c_{u,v})\in\ell^2(\Wnine^+\times\Wnine^+)\), each sum over a product fiber
is finite. The adjoint of \(\mu_\#\) contains all finite-support vectors
in its domain, so that \(\mu_\#\) is closable. The domain of its
closure is
\[
 \Dom(\overline{\mu_\#})=
 \left\{c\in\ell^2(\Wnine^+\times\Wnine^+):
 \sum_{w\in\Wnine^+}
 \left|\sum_{u\boxtimes_\#v=w}c_{u,v}\right|^2<\infty
 \right\}.
\]
Each fiber is finite because $\operatorname{ev}_9$ identifies it with the ordered
pairs of divisors of an integer. The same description holds for
\(\overline{\mu_{\N}}\) after applying \(W_\times\otimes W_\times\), and
$W_\times$ intertwines the graphs of both closures.

\section{Native coproduct and native antipode}

The multiplicative opening is now defined before looking at the codomain:
\[
 \Delta_\#e_w
 =\sum_{u\boxtimes_\#v=w}e_u\otimes e_v.
\]
The sum is obtained by enumerating products of the transducer. One does not ask whether
one ordinary integer divides another.
In the integer sector define, on the corresponding algebraic kernel,
\[
 \Delta_\times e_n
 =\sum_{\substack{a,b\in\N_{\geq1}\\ab=n}}e_a\otimes e_b.
\]
Both coproducts enumerate ordered pairs.

\begin{theorem}[Native coassociativity]
In the algebraic core,
\[
 (\Delta_\#\otimes I)\Delta_\#
 =(I\otimes\Delta_\#)\Delta_\#.
\]
Moreover,
\[
 (W_\times\otimes W_\times)\Delta_\#
 =\Delta_\times W_\times.
\]
\end{theorem}

\begin{proof}
Both sides of the first equality enumerate every native triple
$a\boxtimes_\#b\boxtimes_\#c=w$ exactly once. The second equality follows from the bijection
between fibers induced by \eqref{int:eq:ev-multiplicative}.
\end{proof}

\begin{theorem}[Complete arithmetic coproduct]
\label{int:pf84:C06}
In the algebraic kernel
\(c_{00}(\mathbb N_{\geq1})\subset\ell^2(\mathbb N_{\geq1})\), let
\[
 \Delta_\times e_n=\sum_{ab=n}e_a\otimes e_b.
\]
Its closure has domain
\[
 \operatorname{Dom}(\overline{\Delta_\times})
 =
 \left\{x=(x_n):
 \sum_{n\geq1}d(n)|x_n|^2<\infty\right\},
\]
is coassociative on the kernel and, for \(Qe_n=ne_n\), satisfies
\[
 (\log Q\otimes I+I\otimes\log Q)\Delta_\times
 =\Delta_\times\log Q.
\]
\end{theorem}

\begin{proof}
The supports of \(\Delta_\times e_n\) are orthogonal for distinct integers,
and their squared norm is the number \(d(n)\) of divisors, which yields the weighted
domain of the closure. The two coassociativity relations reindex the same sum
over \(abc=n\). Finally, \(\log a+\log b=\log n\) in each term
proves the intertwining. A triple factorization counted differently, or failure of
this identity, would falsify the theorem.
\end{proof}

\begin{theorem}[Reduced coproduct and irreducible selector]
\label{int:pf84:C07}
In the sector integer, the coproducto reducido
\[
 \Delta_\times^{\rm red}e_n
 =\sum_{\substack{ab=n\\a,b>1}}e_a\otimes e_b
\]
is closable. The operator
\[
 N_\times^{\rm red}
 =\bigl(\overline{\Delta_\times^{\rm red}}\bigr)^*
  \overline{\Delta_\times^{\rm red}}
\]
is positive self-adjoint and diagonal:
\[
 N_\times^{\rm red}e_n=d_{\rm red}(n)e_n.
\]
Its kernel is generated by \(e_1\) and by the \(e_p\) with \(p\) irreducible;
therefore,
\[
 P_{\Irr}
 =\mathbf1_{\{0\}}(N_\times^{\rm red})
  -\lvert e_1\rangle\langle e_1\rvert
\]
select the irreducibles and exclude the unit.

In the native sector, the corresponding coproduct is
\[
 \widetilde\Delta_\#e_w
 =\sum_{\substack{u\boxtimes_\#v=w\\u,v\ne(1)}}e_u\otimes e_v.
\]
The entrelazador \(W_\times\) transports
\(\Delta_\times^{\rm red}\), \(N_\times^{\rm red}\) and \(P_{\Irr}\) toward
\(\widetilde\Delta_\#\),
\[
 A_\#:=
 \bigl(\overline{\widetilde\Delta_\#}\bigr)^*
 \overline{\widetilde\Delta_\#},
 \qquad
 P_{\Irr,\#}
 =(I-P_{(1)})\mathbf1_{\{0\}}(A_\#),
 \qquad
 P_{(1)}=\lvert e_{(1)}\rangle\langle e_{(1)}\rvert.
\]
\end{theorem}

\begin{proof}
The supports of the images of two distinct basis vectors are orthogonal.
The squared norm of
\(\Delta_\times^{\rm red}e_n\) is \(d_{\rm red}(n)\), which vanishes
exactly for \(n=1\) or \(n\) irreducible. This proves the diagonalization
and the declared kernel; subtracting the projector of \(e_1\) removes the unit.
The multiplicativity of \(W_\times\) identifies the integer fibers with the
native fibers and transports the three operators. A composite in the kernel,
an irreducible outside it, or inclusion of the unit in \(P_{\Irr}\)
would refute the theorem.
\end{proof}

\begin{proposition}[Nonadic coincidence of primes and composites]
\label{int:pf84:B05}
Let \(\gcd(n,90)=1\),
\(\tau_n=\operatorname{ord}_n(10)\) and
\[
 M_n=\frac{10^{\tau_n}-1}{n}.
\]
Then \(M_n\equiv0\pmod9\), both for primes and for composites. In
particular, the nonadic closure and the isolated decimal period are not a
selector of irreducibles.
\end{proposition}

\begin{proof}
Because \(9\mid10^{\tau_n}-1\) and because \(n\) is invertible modulo nine, divisibility by nine survives passage to the quotient. The primes \(7,13\) and the composites \(77,91,143\) share period six; the Carmichael numbers provide additional controls. The statement would be falsified only if the congruence or the period separated all primes from all composites in the declared domain.
\end{proof}

This statement is strictly arithmetic: it identifies factorization
fibers and the operators they induce on the sector of normal
forms.

\ifdefined\HMTInternalTraceability
\section{Monoidal extension of the native product to compatible histories}
\label{int:pf84:SYN-05D-HIST}

The result distinguishes the space \(\Krad\) of canonical normal forms and
the space \(\mathcal K_{\#,\mathrm{hist}}\) of complete histories of
cells and propagated quotients.
Many histories may produce the same normal form. Therefore,
$W_\times$ is unitary on $\Krad$, but the projection of a history onto
its value prevents it from being injective on all transport classes.

Let \(\mathscr X_n^{\rm hist}\) be the set of enriched histories of
depth \(n\) and let
\(\mathcal H_n=\ell^2(\mathscr X_n^{\rm hist})\). Denote by \(P_n\) the
orthogonal projection onto the subspace generated by surviving
histories, and suppose a multiplication \(\mu_n\) is defined on an
algebraic domain of
\(\mathcal H_n\otimes_{\mathrm{alg}}\mathcal H_n\). A monoidal extension
of the word product to the inverse system consists of a family of
maps \(\iota_n:\Krad\to\mathcal H_n\) that commutes with the
depth restrictions and satisfies, for every \(n\),
\[
 P_n\,\iota_n\mu_\#
 =P_n\,\mu_n
 (\iota_n\otimes\iota_n)
\]
in the algebraic kernel. This equality is the compatibility condition that
defines the extension. The preceding theorems establish neither the existence nor
the uniqueness of a family \((\iota_n)_n\); they prove the product in the
monoid \(\WordMonoid\) and determine the equation that an extension to the history
space must satisfy. The transfer therefore retains the status of an
\emph{open program}: it does not inherit the exactness of the product in
\(\WordMonoid\). The nonexistence of a compatible family, failure to
commute with truncations, loss of survival, quotient,
carry, orientation, or memory, or failure to satisfy the preceding identity
are its falsifiers.
\fi

\section{Exhaustive enumeration and necessity of structural coordinates}

The finite enumeration comprises:
\[
\begin{array}{@{}lr@{}}
\toprule
\text{Test}&\text{Result}\\\midrule
\text{Cells per local map}&81\\
\text{Reversible numbering paths}&59050\\
\text{Intertwiner pairs}&532900\\
\text{Intertwiner errors}&0\\
\text{Associative triples}&531441\\
\text{Associative errors}&0\\
\text{Coproduct window}&1\text{ to }500\\
\text{Irreducibles obtained}&95\\
\bottomrule
\end{array}
\]

\begin{definition}[Multiplicative quotient]
Preserving only $R^\times$ transforms $2\cdot5=10$ into the visible value $1$.
Perturbing $q^\times(9,9)$ from $8$ to $7$ transforms $81$ into $72$. The local
quotient is necessary.
\end{definition}

\begin{definition}[Additive application]
The multiplicative map produces partial products, but does not resolve their collisions.
Removing $q^+$ breaks the global normalization; for example, it no longer reproduces
the intermediate quotients of $10\cdot18=180$.
\end{definition}

\begin{definition}[Coordinates and boundary symbol]
Ordinary zero is not a digit of the construction. The empty word is
$0$, and the symbol $9$ preserves its value and quotient. Mixing
zero-indexed positional coordinates with the labels $1,\ldots,9$ changes the
cells.
\end{definition}

\begin{proposition}[Exact scope of the ablations and of the finite census]
\label{int:pf84:D10}
The preceding ablations prove, within the declared conventions:
\begin{enumerate}
  \item without \(q^\times\), \(2\cdot5=10\) colapsa to the residue visible \(1\);
  \item cambiar \(q^\times(9,9)=8\) by \(7\) transforms \(81\) in \(72\);
  \item without \(q^+\), normalization of the additive collisions fails.
\end{enumerate}
The exhaustive enumeration comprises the \(2\cdot81=162\) cells of the two
maps, the \(729\) local triples, \(532\,900\) word pairs, and
the \(95\) irreducibles no greater than \(500\). The indicated identities are
proved on these finite domains, but the coinductive survival
of every TPK history is not.
\end{proposition}

\begin{proof}
The three counterexamples are computed directly with
\eqref{int:eq:local-plus} and \eqref{int:eq:local-times}. The enumeration of the four
finite domains produces the cardinals written in the statement and does not add
any coinductive identity to those already proved in the monoid
of words.
\end{proof}

\section{Algebraic structure of the word monoid}

The construction proves, in the monoid \(\WordMonoid\):
\begin{enumerate}
 \item a unique normal form for each nonnegative integer;
 \item global sum and product constructed solely from the two local maps;
 \item symbolic conservation of the value under quotient propagation;
 \item a unitary $W_\times$ and the exact multiplicative square;
 \item a divisor coproduct and an irreducible selector defined natively;
 \item the need to preserve the transport class when leaving the
normal-form sector.
\end{enumerate}

Passage to the enriched space now reduces to an existence problem:
constructing a family \((\iota_n)_n\) compatible with the depth restrictions
and proving that it satisfies the equation of the preceding section.
This question is distinct from the exact intertwiner on
\(\WordMonoid\), which is already determined by the normal form and the
local maps.



\label{p3-83-c23-s1:ch:foundation}

\section{Problem of determining the constants from the genealogical state}

A physical constant is represented as a section of a magnitude line equipped with a rule of evaluation, orientation, and normalization. Its numerical coordinate depends on the chosen chart, whereas its genealogy remains invariant. Consequently, the determination problem comprises:

\begin{enumerate}
\item the identification of the antecedent HMT state;
\item the specification of the operator or character of selection;
\item determination of the conserved invariant;
\item the corresponding magnitude line;
\item the chart in which the terminal value is expressed.
\end{enumerate}

The causal chain of the work is

\begin{equation}
\begin{aligned}
\APP&\longrightarrow\TRIT\longrightarrow\TPK
\longrightarrow\text{internal realization}
\longrightarrow\text{dimensionless invariant}\\
&\longrightarrow\text{magnitude section}
\longrightarrow\R.
\end{aligned}
\label{p3-83-c23-s1:eq:master-chain}
\end{equation}

The term \(\R\) in the chain represents the terminal evaluation of the
section; it does not constitute the original domain from which the
route is retrospectively reconstructed.

\section{Antecedent chain APP--TRIT--TPK and domain of the extraction operators}

APP supplies the geometry of positions, residues, quotients, and
boundaries. TRIT incorporates orientation \((-1,0,+1)\) and distinguishes
opening, threshold, and return. TPK completes the state through carry, route,
sheet, memory, and survival. In particular, transition \(9\to1\)
preserves phase and modifies the memory state; it therefore does not constitute
a reinitialization devoid of genealogy.

The realizations developed below preserve that antecedent without replacing it. The vacuum is obtained through a spectral function of the bilayer operator; gravity selects a section through the chronological structure \(108\); criticality uses a dodecaphasic family; the Gaussian constants follow from the evaluation of cylinders; and the modular Hamiltonian is defined on a boundary that preserves the incidence \(6\to8\). The production of scalar values does not establish functional identity among these realizations.

\begin{proposition}[Insufficiency of the Scalar Coincidence]
Let \(F_1:X_1\to L_1\) and \(F_2:X_2\to L_2\) be two typed functionals. Equality of
coordinates \(c_1(F_1(x_1))=c_2(F_2(x_2))\) does not by itself induce either an
identification \(X_1\simeq X_2\) or an intertwiner
\(F_2I=JF_1\).
\end{proposition}

In particular, a numerical proximity such as
\(I_{\rm or}\approx\gammaH a_0\) cannot be promoted to an identity without
constructing the corresponding intertwiner.

\section{Lines of magnitude, oriented letters and dimensional coordinate systems}

Let \(L_Q\) be an oriented real line representing the magnitude \(Q\).
A unit \(u_Q\in L_Q\setminus\{0\}\) determines a
provisional coordinate \(L_Q\simeq\R\), without replacing the intrinsic structure of
the line. The constitutive relation

\[
L_\mu=L_S L_C^{-1}L_Q^{-2},\qquad
L_\varepsilon=L_S^{-1}L_C^{-1}L_Q^2
\]

precedes the SI choice. In the same way,

\[
L_Z=L_SL_Q^{-2},\qquad L_G=L_C^5L_T^2L_S^{-1}
\]

type impedance and the gravitational constant before expressing their magnitudes in
ohms or SI units. In this work we shall write

\[
Q_{\rm int}=\widehat Q\,u_Q
\]

to distinguish the dimensionless character \(\widehat Q\) from the reference
section \(u_Q\).

\section{Torsor of frames and covariance of dimensional evaluations}

The admissible charts form a torsor and therefore none has a character that is canonical relative to the others. The frame change \(u_Q\mapsto\lambda_Q u_Q\) induces the contragredient transformation of the coordinate. Physical identities must be homogeneous with respect to this action. For example,

\[
\mu\varepsilon=c^{-2},\qquad \frac\mu\varepsilon=Z^2
\]

are valid in every compatible frame. Mathematical unification and
dimensional unification do not eliminate units; rather, they reconstruct them through the
covariance of the lines and allow the invariant to precede its coordinate
expression.

\section{Final metrological evaluation and prohibition of retrospective selection}

Each entry of the final atlas adopts the form

\begin{equation}
\boxed{\text{antecedent}\to\text{operator}\to\text{invariant}
\to\text{meaning}\to\text{terminal value}.}
\label{p3-83-c23-s1:eq:atlas-rule}
\end{equation}

This order expresses a causal dependency: the numerical value is only
determined after fixing the antecedent, operator, invariant, and
interpretation. Omission of any of these elements reduces the
result to a reconstruction that is numerical, not a derivation that is
genealogical.

\section{Exact domain of the genealogical generation and its terminal metrical chart}

The generation of
\(\pi,e,\varphi\), the dodecaphasic closure of \(\alpha_{\HMT}\),
the action section, and the Barbero--Immirzi output \(\gammaH\) are adopted as initial results. The
present development does not recalibrate those constructions; it determines the constants
and invariants resulting from their application to the constitutive,
thermal, gravitational, modular, and spectral realizations.

The electron mass is likewise not included as a governing problem, since it
belongs to the General Mass Law. Every occurrence of \(m_e\) in an
atomic identity is treated as a datum proceeding from that domain and not as
a new selection of the present development.

\subsection*{Logical Scope and Domain of Validity}
The APP--TRIT--TPK architecture and the dimensional lines constitute the
starting structure. The uniform rule of five fields of the atlas makes
explicit its causal discipline in each entry.



\section{Functional taxonomy of \(114\) numerical expressions and \(106\) operator readings}
\label{p3-83-c23-s2:sec:atlas-taxonomico-114-constantes-rev11}
\index{constants!taxonomic atlas of 114 entries}

Let the finite indexed set be
\[
 E_{114}=\{\varepsilon_1,\ldots,\varepsilon_{114}\}
\]
of expresiones numericas and let \(\mathcal O\) the space of readings of the
reader triadico. The map
\[
 T:E_{114}\longrightarrow\mathcal O,
 \qquad \varepsilon_i\longmapsto T_{\rm tri}(\varepsilon_i),
 \qquad |\operatorname{im}T|=106,
\]
classifies expressions according to their operatorial role. Each
\(\varepsilon_i\), including the compared targets, is fixed as
input of \(T\); therefore, this map does not generate the values that
it evaluates. The functional partition is
\[
\begin{array}{l|r}
\text{functional family}&\text{entries}\\ \hline
\text{structural kernel}&5\\
\text{coupling and closure}&1\\
\text{bridges between modes}&7\\
\text{angular and closure}&22\\
\text{exponential--logarithmic}&12\\
\text{spectral and regularization}&14\\
\text{hyperbolic and modular}&2\\
\text{algebraic and metric}&50\\
\text{typed miscellany}&1
\end{array}
\qquad
\sum n_i=114.
\]
Network centrality, shortest paths, and communities computed on
the graph induced by \(T\) describe connectivity within the classification;
they do not establish mathematical causality or HMT genealogy.

Seven fibers of \(T\) gather fifteen expressions with a common reading. A
shared fiber does not identify their values either. In particular,
\(\varphi^{-1}\), \(\varphi\), and \(\varphi^2\) are distinct real numbers that the
map \(T\) groups into the same self-scaling orbit; the common
projection preserves dynamic function and forgets the Archimedean coordinate that
separates them.

The exhaustive enumeration traversed the \(114\) entries of \(E_{114}\), four hundred sixty-eight direct candidates \(U_{12}\), and one hundred four thousand nine hundred seventy-six weighted seeds without finding, within that finite domain, a rival that reproduced the oriented dodecaphasic correction of \(\alpha_{\rm HMT}\). The search domain is thereby bounded by those three sets. The census proves that no rival appears in the enumerated class; it does not prove uniqueness among all possible functionals and does not enter as a premise of theorem \(\alpha_{\rm HMT}\).

For each \(\varepsilon_i\), the taxonomy assigns the quadruple
\((\text{object},\text{reader},\text{domain},\text{structural function})\)
used in the enumeration.
